\section{Matrix metrics and flow comparison}\label{sec:matrix}

\subsection{The finite flow}
The space $\cP_r$ is a complete simply connected manifold of
nonpositive sectional curvature. The geodesic from $P$ to $Q$ is
\begin{equation}\label{eq:geodesic}
 \Geo(P,Q;\theta)=P^{1/2}
   \exp\!\bigl(\theta\log(P^{-1/2}QP^{-1/2})\bigr)P^{1/2}.
\end{equation}
Its initial velocity is
$\Log_PQ=P\log(P^{-1}Q)$. These formulas and the distance
\eqref{eq:target-metric} are standard \cite[Sections~2.3--2.6]{bhatia}.
The determinant-one space is totally geodesic: the logarithm of the
determinant along \eqref{eq:geodesic} is the affine interpolation of
the endpoint logarithmic determinants. Squared distance is jointly
convex along pairs of geodesics, and every two points have a unique
geodesic. All maps in \eqref{eq:geodesic} are smooth, including at
repeated eigenvalues.

\begin{lemma}\label{lem:finite-flow}
Equation~\eqref{eq:finite-flow} is the negative gradient flow of
\eqref{eq:graph-energy} for the product metric with weights $A_i$.
It preserves determinant one and exists for all $\tau\ge0$ in a
compact subset depending only on its initial tuple and the areas.
All its derivatives of each fixed order are bounded on bounded
time intervals.
\end{lemma}
\begin{proof}
The gradient of $\frac12d(P,Q)^2$ in the first variable is
$-\Log_PQ$; applying this identity at each vertex gives
\eqref{eq:finite-flow}, including the factor $2\pi$ and the weight
$A_i^{-1}$. The same intrinsic graph-gradient formula appears in
\cite[Section~III]{tav}.
Since $\tr\log(H_i^{-1}H_j)=0$, determinant one is preserved.

Let $o=(I,\ldots,I)$ in the product, and write $d_A$ for its weighted
distance. Geodesic convexity and $\mathscr E(o)=0$ give
\[
 \langle\nabla\mathscr E(\mathbf H),\Log_{\mathbf H}o\rangle_A
 \le-\mathscr E(\mathbf H).
\]
Along the negative gradient flow this is
\[
 \frac d{d\tau}\frac12d_A(\mathbf H,o)^2
 \le-\mathscr E(\mathbf H)\le0.
\]
Closed bounded sets in this finite-dimensional complete manifold are
compact. The smooth vector field therefore continues for all time,
and differentiation of the ordinary differential equation gives the
stated derivative bounds.
\end{proof}

\subsection{The exact HYM equation}
In a holomorphic coordinate $z=x+iy$, write
$\omega_t=(i/2)a(z)\,dz\wedge d\bar z$. The nonnegative real
scalar Laplacian is $\Delta_t=-4a^{-1}\partial_z\partial_{\bar z}$.
Equation~\eqref{eq:hym-intro} reads
\begin{equation}\label{eq:complex-hym}
 \partial_sH=4a^{-1}(H_{z\bar z}-H_{\bar z}H^{-1}H_z).
\end{equation}
Indeed $F_H=-\partial_{\bar z}(H^{-1}H_z)\,dz\wedge d\bar z$.
Expanding the complex derivatives gives
\begin{equation}\label{eq:real-hym}
 \begin{split}
 a\partial_sH={}&H_{xx}+H_{yy}-H_xH^{-1}H_x-H_yH^{-1}H_y\\
 &+i(H_xH^{-1}H_y-H_yH^{-1}H_x).
 \end{split}
\end{equation}
The final term is Hermitian. It is generally nonzero and is absent
from ordinary symmetric-space harmonic-map heat flow. For example,
take $H(x,y)=\exp(xU+yV)$ with noncommuting trace-free Hermitian
matrices $U,V$. At the origin the first four terms cancel, whereas
the last is $i[U,V]$. The proof below makes no assumption that this
term vanishes for an actual solution.

In rank one $H=e^u$ gives $\partial_su=-\Delta_tu$. Taking the
trace of the logarithmic velocity in general gives the same heat
equation for $\log\det H$, so the determinant normalization persists.
Linearization at a constant identity matrix is consequently
$\partial_su=-\Delta_tu$ on trace-free Hermitian entries.

\subsection{A comparison estimate with a residual}
For a smooth comparison family $K(s)$ define
\begin{equation}\label{eq:residual}
 R(K)=K^{-1}\partial_sK+2i\Lambda_{\omega_t}F_K.
\end{equation}
It is self-adjoint with respect to $K$. Its pointwise norm below is
the operator norm in that metric.

\begin{lemma}\label{lem:comparison}
Let $H(s)$ be an exact HYM solution on a fixed compact fibre, and
$K(s)$ any smooth positive comparison family. With $h=H^{-1}K$,
\begin{equation}\label{eq:forced-sigma}
 (\partial_s+\Delta_t)\sigma(H,K)
 \le\tr((h-h^{-1})R(K)).
\end{equation}
If $mI\le H,K\le MI$ throughout the interval, then
\begin{equation}\label{eq:integrated-sigma}
 \int_{X_t}\sigma(H(s),K(s))\,dA_t
 \le\int_{X_t}\sigma(H(0),K(0))\,dA_t
 +\frac{2rM}{m}\int_0^s\!\int_{X_t}|R(K)|\,dA_t\,du.
\end{equation}
The constants do not use a spectral gap or a Sobolev constant of $X_t$.
\end{lemma}
\begin{proof}
For two exact flows, the Donaldson trace calculation gives
$(\partial_s+\Delta_t)\sigma\le0$; see also
\cite[Definition~2.3 and Proposition~2.4]{wang}, with the bundle
morphism zero. We recall the pointwise calculation to fix conventions.
Choose a time-independent holomorphic frame normal for $H$ at the
point, with $H=I$ and $\partial H=0$. The curvature terms cancel
between the time derivative and the scalar Laplacian of
$\tr(H^{-1}K)$, leaving
\[
 -4a^{-1}\tr(K_{\bar z}K^{-1}K_z)\le0.
\]
This is nonpositive because $K_{\bar z}=K_z^*$. Interchanging
$H,K$ proves the same assertion for the inverse trace.

Relative to the exact comparison velocity, the actual velocity of
$K$ has the additional term $KR(K)$. Its contribution to the two
trace derivatives is respectively $\tr(hR)$ and $-\tr(Rh^{-1})$.
This proves \eqref{eq:forced-sigma}. Both $h$ and $R$ are
$K$-self-adjoint, since $Kh=KH^{-1}K$ is Hermitian. Therefore
\[
 |\tr((h-h^{-1})R)|
 \le(\tr h+\tr h^{-1})|R|
 \le 2r(M/m)|R|.
\]
Integration over the whole compact fibre removes the Laplacian;
the metric and area form are fixed during the flow. This gives
\eqref{eq:integrated-sigma}.
\end{proof}

If $e^{b_j}$ are the eigenvalues of $H^{-1}K$, then
\[
 d(H,K)^2=\sum_j b_j^2
 \le\sum_j(e^{b_j}+e^{-b_j}-2)=\sigma(H,K),
\]
which proves the inequality in \eqref{eq:sigma}.

\begin{lemma}\label{lem:existence-bounds}
For every fixed smooth compact fibre and smooth positive initial
metric, the HYM metric flow exists smoothly for all $s\ge0$.
For the trivial holomorphic bundle, initial bounds $cI\le H(0)\le CI$
imply the same bounds at every later time. These comparison constants
are independent of the degeneration parameter.
\end{lemma}
\begin{proof}
The first assertion is the classical compact-base existence theorem
for Donaldson's heat flow \cite{donaldson}; the convention
\eqref{eq:hym-intro} and the all-time statement are also given in
\cite[Section~4.1, equation~(4.8) and following paragraph]{hkkp}.
The order-preservation principle
\cite[Proposition~4.3]{hkkp} compares the solution with the two
stationary constant matrix metrics $cI,CI$. Neither comparison
changes as the base parameter varies. Determinant one follows from
the scalar heat equation for $\log\det H$.
\end{proof}
